\item \textbf{Выбор уставки токового органа пуска охлаждения}

%\begin{enumerate}[label=\arabic{section}.\arabic{enumi}.\arabic*, labelsep=4pt, leftmargin=0em, itemindent=65pt, parsep=3pt]
\begin{enumerate}[label={}, labelsep=4pt, leftmargin=0pt, itemindent=57pt]

\item
По выражению (\ref{eq:rtpo}) определяется первичный ток органа пуска охлаждения $I_{\textsl{ср}}$:
\begin{equation*}
    I_{\textsl{ср}} = K_{\textsl{отс}}\cdot I_{\textsl{ном.ст.}}, = 0,95 \cdot 250 = 238\; (\textsl{А}),
\end{equation*}

Уставка задается в относительных единицах от номинального тока аналогового входа устройства, для этого полученное значение пересчитывается по следующему выражению:
\begin{equation*}
    I_{\textsl{ср}} = \frac{I_{\textsl{с.з.ЗП}}}{I_{\textsl{ном.вх.уст}}} \cdot \frac{I_{\textsl{ном.вт.ТТ}}}{I_{\textsl{ном.пер.ТТ}}} = \frac{238}{5} \cdot \frac{5}{400} = 0,6\; (\textsl{о.е.}),
\end{equation*}
\begin{eqexpl}[15mm]
    \item{$I_{\textsl{ном.вх.уст}}$} номинальный ток аналоговых входов устройства, 1 или 5 А;
    \item{$I_{\textsl{ном.вт.ТТ}}$} номинальный вторичный ток ТТ, А;
    \item{$I_{\textsl{ном.пер.ТТ}}$} номинальный первичный ток ТТ, А.
\end{eqexpl}

Параметр \textbf{$I_{\textsl{ср}}$} <<Ток срабатывания>> функции РТПО принимается \textbf{равным 0,6}.

\end{enumerate}
